\section{Conclusión individual: César Macías}

Mi responsabilidad en este proyecto abarcó el desarrollo del cliente gRPC
(\texttt{src/grpc/grpc\_client.py}), la ejecución y automatización de la batería
de experimentos distribuidos (\texttt{src/grpc/grpc\_experiment.py}) y la
implementación de las pruebas unitarias correspondientes en
\texttt{tests/test\_grpc\_client.py}.

La sincronización de los relojes lógicos de Lamport demostró ser fundamental para
mantener la consistencia temporal en un entorno asíncrono y distribuido. Implementar
la actualización del reloj antes de emitir la llamada y tras la llegada del acuse
permitió constatar que la relación de precedencia causal se cumple con exactitud en
todas las corridas, tanto en los intercambios secuenciales como bajo condiciones
concurrentes. Asimismo, la serialización binaria con Protocol Buffers redujo la
complejidad del cliente respecto a TCP, al delegar la codificación de tipos y
ofrecer stubs tipados que validan la interfaz antes del tránsito por la red.

Las pruebas unitarias implementadas verificaron no solo el caso nominal, sino
también la robustez ante destinos caídos y validaciones de argumentos. Con estas
adiciones, la suite de pruebas del proyecto alcanzó 56 pruebas unitarias exitosas,
garantizando la fiabilidad de las mediciones y la base de datos experimental
utilizada para la comparación estadística entre ambos paradigmas.
